\chapter{RESPONSABILIDADE SOCIAL, DIVERSIDADE E INCLUSÃO}

Os usuários do aplicativo são, em grande parte, profissionais do armazém, da portaria e da frota, funções historicamente ocupadas por pessoas negras e de baixa renda e por trabalhadores com pouca familiaridade digital. Por isso, a solução foi pensada para não criar barreiras e para reduzir as existentes, em linha com o Estatuto da Igualdade Racial \cite{brasil2010igualdade} e com a Lei Brasileira de Inclusão \cite{brasil2015lbi}.

\section{Combate à discriminação e valorização da diversidade}

O sistema registra fatos operacionais (quem bipou, quando e o quê), e não avaliações subjetivas sobre as pessoas. Assim, uma divergência na conferência é tratada como dado registrado, e não como acusação verbal. Os perfis de acesso são definidos pela função, nunca por características pessoais.

Para valorizar a diversidade, propõe-se acompanhar a implantação de um programa de formação digital para as equipes operacionais, com prioridade à promoção interna de conferentes e operadores. Nas datas de 20 de novembro (Consciência Negra) e 19 de abril (Povos Indígenas), a empresa pode promover rodas de conversa sobre a história e a cultura afro-brasileira e indígena, tema que a Lei n.º 11.645 \cite{brasil2008} tornou obrigatório na educação básica.

\section{Recursos tecnológicos inclusivos e acessibilidade}

O aplicativo foi desenhado seguindo princípios de design de interação centrado no usuário \cite{preece2013} e as diretrizes de acessibilidade WCAG 2.1 \cite{w3c2018}. Os botões têm 64 pixels de altura e o texto de 16 a 18 pontos, para uso com uma só mão ou com luvas. Sucesso e erro geram vibrações diferentes, para que o operador perceba o resultado sem olhar a tela, e as cores vêm sempre acompanhadas de texto e ícone, para não depender da percepção de cor. Os indicadores de progresso têm rótulo para leitores de tela, o aplicativo respeita o tamanho de fonte e o tema do sistema, e o código sempre pode ser digitado em vez de lido pela câmera. As mensagens de erro usam linguagem simples e dizem o motivo da recusa, como ``faltam 8 caixas para bipar''.

\section{Tecnologia para a promoção da cidadania}

A proteção de dados pessoais também é uma questão de cidadania. Em conformidade com a LGPD \cite{brasil2018lgpd}, o CPF dos motoristas aparece mascarado nas telas, as senhas são guardadas apenas como \textit{hash} e cada perfil acessa somente os dados necessários à sua função. Os registros objetivos também protegem o trabalhador: o motorista que registra uma não entrega com motivo e horário tem prova de que cumpriu a rota.
